\documentclass[../DoAn.tex]{subfiles}
\begin{document}

Chương 2 cung cấp bức tranh toàn cảnh về cơ sở lý thuyết và các nghiên cứu nền tảng liên quan đến đề tài. Đầu tiên, chương sẽ làm rõ khái niệm về hệ thống phát hiện xâm nhập mạng (NIDS) cùng kiến trúc lai (Hybrid IDS). Tiếp theo, điểm qua các nghiên cứu tương tự đã được thực hiện trong lĩnh vực này, từ đó rút ra ưu nhược điểm làm tiền đề cho đề xuất của đồ án. Cuối cùng, chương đi sâu vào nền tảng của các thuật toán Học sâu (Deep Learning) cũng như các công nghệ mã nguồn mở đóng vai trò then chốt trong việc xây dựng hệ thống.

\section{Hệ thống phát hiện xâm nhập mạng (NIDS) và Kiến trúc Hybrid IDS}
Hệ thống phát hiện xâm nhập mạng (NIDS - Network Intrusion Detection System) là một thành phần không thể thiếu trong cấu trúc an ninh mạng hiện đại. Khác với tường lửa (Firewall) vốn thường chỉ chặn các kết nối dựa trên quy tắc tĩnh, NIDS có khả năng phân tích sâu nội dung gói tin để tìm ra dấu hiệu của mã độc hoặc hành vi bất thường trong toàn bộ mạng lưới.

Theo phương pháp phân tích, NIDS được chia làm hai loại chính. Loại thứ nhất là phân tích dựa trên dấu hiệu (Signature-based), so khớp nội dung gói tin với một cơ sở dữ liệu các mẫu tấn công đã biết trước. Loại thứ hai là phân tích dựa trên bất thường (Anomaly-based), thiết lập "đường cơ sở" (baseline) của hoạt động mạng bình thường và đánh dấu sự sai lệch là một mối đe dọa. Mỗi phương pháp đều có điểm mạnh riêng: Signature-based hoạt động nhanh và chính xác với các lỗi đã biết nhưng không thể đối phó với tấn công Zero-day, trong khi Anomaly-based có thể phát hiện Zero-day nhưng thường xuyên gây ra cảnh báo giả.

Kiến trúc Hybrid IDS (Kiến trúc lai) ra đời nhằm dung hòa hai phương pháp trên. Bằng cách sử dụng một hệ thống NIDS Signature-based ở vòng ngoài để lọc và cảnh báo tức thì các cuộc tấn công phổ biến, hệ thống giảm tải đáng kể khối lượng tính toán. Phần lưu lượng phức tạp tiếp tục được đưa qua module Anomaly-based (thường tích hợp AI) ở vòng trong để phân tích sâu hơn. Sự kết hợp này mang lại hệ thống phòng thủ toàn diện với độ chính xác cao và độ trễ thấp.

\begin{table}[htbp]
\centering
\caption{So sánh giữa Signature-based IDS và Anomaly-based IDS}
\label{tab:ids_comparison}
\begin{tabular}{|p{3.5cm}|p{5.5cm}|p{5.5cm}|}
\hline
\textbf{Tiêu chí} & \textbf{Signature-based IDS} & \textbf{Anomaly-based IDS} \\
\hline
\textbf{Cơ chế hoạt động} & So khớp mẫu (pattern) với CSDL đã biết & Phân tích hành vi, đối chiếu đường cơ sở \\
\hline
\textbf{Phát hiện Zero-day} & Không có khả năng & Có khả năng nhận diện \\
\hline
\textbf{Cảnh báo giả (FPR)} & Rất thấp & Thường ở mức cao \\
\hline
\textbf{Bảo trì} & Cao (Cần cập nhật luật thường xuyên) & Trung bình (Cần huấn luyện lại mô hình) \\
\hline
\textbf{Tốc độ \& Tài nguyên} & Rất nhanh, ít tốn kém & Chậm hơn, tốn tài nguyên tính toán \\
\hline
\end{tabular}
\end{table}

\section{Các kết quả nghiên cứu tương tự về ứng dụng AI trong NIDS}
Trong thập kỷ qua, sự bùng nổ của Trí tuệ nhân tạo (AI) đã thúc đẩy mạnh mẽ các nghiên cứu ứng dụng Học máy và Học sâu vào NIDS. Nhiều công trình nghiên cứu sử dụng các bộ dữ liệu tiêu chuẩn như KDD Cup 99, NSL-KDD hay mới hơn là CIC-IDS-2017 để huấn luyện mô hình.

Các thuật toán Học máy cổ điển như Cây quyết định (Decision Tree), Rừng ngẫu nhiên (Random Forest) hay K-Nearest Neighbors (KNN) đã chứng minh được hiệu quả trong việc phân loại tấn công. Tuy nhiên, các thuật toán này thường gặp khó khăn trước các cuộc tấn công tinh vi hoặc khi phải đối mặt với bài toán dữ liệu mất cân bằng nghiêm trọng (imbalanced data).

Để khắc phục, các mạng nơ-ron học sâu như Mạng nơ-ron nhân tạo (ANN), Mạng tích chập (CNN) và Mạng bộ nhớ dài-ngắn (LSTM) bắt đầu được áp dụng. Dù có khả năng tự động trích xuất đặc trưng ẩn sâu, kiến trúc của chúng khá cồng kềnh, tiêu tốn tài nguyên và có thời gian suy luận (inference time) lớn. Gần đây, sự xuất hiện của mô hình Transformer đã mở ra hướng đi mới, đặc biệt là các biến thể xử lý dữ liệu dạng bảng, cung cấp khả năng hiểu sâu mối tương quan của các đặc trưng mạng với chi phí tính toán hợp lý.

\section{Cơ sở lý thuyết về Học sâu: Autoencoder, FT-Transformer và Ensemble Learning}
Trong đồ án này, ba kỹ thuật tiên tiến được sử dụng làm nền tảng lý thuyết AI cốt lõi:

\textbf{Autoencoder}: Là một mạng nơ-ron học không giám sát (unsupervised learning) được thiết kế để nén dữ liệu đầu vào thành một biểu diễn ẩn và sau đó tái tạo lại nó. Trong NIDS, mạng được huấn luyện bằng dữ liệu mạng bình thường; do đó, khi một gói tin tấn công đi qua, lỗi tái tạo (reconstruction error) sẽ tăng đột biến, giúp nhận diện sơ bộ sự bất thường để lọc bớt lưu lượng an toàn.

\begin{figure}[htbp]
    \centering
    \includegraphics[width=0.7\textwidth]{Hinh_ve/autoencoder_arch.png}
    \caption{Cấu trúc cơ bản của mạng Autoencoder trong nhận diện bất thường}
    \label{fig:autoencoder}
\end{figure}

\textbf{FT-Transformer (Feature Tokenizer Transformer)}: Là biến thể của mô hình Transformer nổi tiếng, tối ưu hóa cho dữ liệu dạng bảng (tabular data). Khác với mạng ANN truyền thống, FT-Transformer biến đổi từng đặc trưng (feature) thành một token riêng biệt, sau đó áp dụng cơ chế tự chú ý (Self-Attention) để nắm bắt mối quan hệ phức tạp giữa tất cả các đặc trưng. Điều này giúp mô hình trích xuất thông tin có ý nghĩa nhất từ vectơ dữ liệu nhiều chiều.

\begin{figure}[htbp]
    \centering
    \includegraphics[width=0.8\textwidth]{Hinh_ve/ft_transformer_arch.png}
    \caption{Cấu trúc mạng FT-Transformer tối ưu hóa cho dữ liệu dạng bảng}
    \label{fig:ft_transformer}
\end{figure}

\textbf{Ensemble Learning (Học quần thể)}: Là phương pháp kết hợp dự đoán của nhiều mô hình đơn lẻ (như FT-Transformer, Random Forest, KNN) thông qua cơ chế biểu quyết (Voting) nhằm đưa ra kết quả cuối cùng ổn định hơn. Học quần thể giúp khắc phục triệt để điểm yếu cục bộ của từng thuật toán, đặc biệt là giảm tỷ lệ cảnh báo giả xuống mức cực thấp đối với các dạng tấn công phức tạp.

\section{Công nghệ triển khai: Snort, CICFlowMeter và Kiến trúc Streaming}
Để biến lý thuyết thành một hệ thống thực chiến, đồ án sử dụng kết hợp nhiều công nghệ hiện đại:

\textbf{Snort}: Là một NIDS mã nguồn mở hoạt động dựa trên luật (Signature-based) phổ biến nhất thế giới. Snort thực hiện việc kiểm tra gói tin theo thời gian thực và đóng vai trò như "lưới lọc" tầng đầu tiên trong kiến trúc Hybrid của đồ án.

\textbf{CICFlowMeter}: Là công cụ trích xuất đặc trưng lưu lượng mạng có khả năng nhận các luồng gói tin và tự động tính toán ra hàng chục đặc trưng thống kê (như tỷ lệ luồng, số lượng gói tin lỗi). Dữ liệu này là đầu vào chuẩn hóa hoàn hảo cho các mô hình AI.

\textbf{Kiến trúc Streaming và Web API}: Dữ liệu sau khi trích xuất được đưa vào luồng dữ liệu thời gian thực. Đồ án ứng dụng FastAPI để đóng gói mô hình AI thành một microservice hiệu năng cao, nhận dữ liệu liên tục và trả về nhãn phân loại trong vòng vài mili-giây. Kết quả cảnh báo sau đó được đẩy lên Dashboard giám sát bằng React và Node.js.

\vspace{1em}
\textbf{Kết chương}
Chương 2 đã phác họa bức tranh toàn cảnh về lý thuyết và công nghệ đằng sau hệ thống phát hiện xâm nhập mạng. Thông qua việc phân tích điểm hạn chế của các nghiên cứu cũ và làm rõ cơ chế của các kỹ thuật tiên tiến, chúng ta đã có cơ sở vững chắc để thiết kế kiến trúc hệ thống Hybrid NIDS mới. Những ý tưởng nền tảng này sẽ được hiện thực hóa chi tiết trong Chương 3.

\end{document}
